\section{Related Works}
\label{sec:relatedworks}

\subsection{Video Large Language Models}

With the rapid progress of Large Language Models (LMMs)~\cite{achiam2023gpt,chiang2023vicuna,taori2023stanford} and Multimodal Large Language Models (MLLMs)~\cite{alayrac2022flamingo,li2023blip,liu2024improved,liu2024llavanext,team2023gemini,bai2025qwen2,zhu2025internvl3,zhang2025omnicharacter,zhang2025text}, nowadays there has been growing interest in developing Video LLMs (VLLMs)~\cite{bai2025qwen2,zhu2025internvl3,zhang2023video,lin2023video,cheng2024videollama,li2024mvbench,weng2024longvlm,wu2023visual,wang2023chatvideo,li2023videochat,zeng2022video,chen2024sharegpt4video}, enabling video understanding and question answer tasks. They can be categorized into general Video LLMs and Video LLMs with visual token compression during training time. General Video LLMs~\cite{lin2023video,li2024llava,an2025llava,zhang2024video,cheng2024videollama,bai2025qwen2,zhu2025internvl3} extract directly the raw video frames tokens and then simply apply pooling before feeding to the LLM. Moreover, to handle spatiotemporal information for videos, some methods~\cite{bai2025qwen2,zhang2024video} extend the 2D image position encoding into video by inducing a temporal dimension. Video LLMs with visual token compression during the training-time~\cite{xu2024pllava,maaz2023video,jin2024chat,shen2024longvu,zohar2025apollo,li2024llama} propose to reduce video tokens significantly to enable long-context video processing. 



However, due to the high spatiotemporal demands of complex video understanding tasks, the visual tokens dominate the following LLM computation overhead since the tokens length can reach up to one million for hours-long videos~\cite{team2024gemini}, substantially increasing inference time and memory consumption. Though some methods~\cite{lin2024vila,liu2025nvila} aim to optimize token utility, they still require model fine-tuning and demand considerable hardware resources. This underscores a critical need for developing more efficient, training-free token compression methods specifically for video LLMs, bypassing the need for costly model adaptations and significant hardware consumption.

\subsection{Image Visual Token Compression}

Token compression has emerged as an effective approach to reduce the computational overhead and complexity in transformer visual encoders and large language models, such as ViT~\cite{dosovitskiy2020image}, CLIP~\cite{radford2021learning} and SigLip~\cite{zhai2023sigmoid}. Pioneering works such as ToMe~\cite{bolya2022token} and FastV~\cite{chen2024image} have explored methods for visual token merging spatially and text-guided pruning to improve the efficiency of LVLMs~\cite{kong2025token}. Hence, this line of approaches can be broadly divided into two main categories: (1) Text-agnostic token compression approaches~\cite{yang2025topv,arif2025hired,shang2025llava,wen2025stop,yang2025visionzip,zhang2024cls,zhang2025vscan}, which discover and merge or remove redundant or uninformative visual tokens during visual encoding. VisionZip~\cite{yang2025visionzip} picks out the dominant visual tokens based on \texttt{[CLS]} attention scores, while LLaVA-PruMerge~\cite{shang2025llava} identifies significant spatial redundancy among visual tokens and applies an adaptive visual token reduction strategy to reduce the number of visual tokens. (2) Text-guided pruning approaches~\cite{tan2025tokencarve,zhang2024sparsevlm,xing2025conical,xing2024pyramiddrop,liu2024multi,ye2025atp,khaki2025sparsevila}, which target at removing visual tokens that are irrelevant to the text query during the LLM decoding phase. SparseVLM~\cite{zhang2024sparsevlm} utilizes an iterative sparsification strategy to exploit visual-relevant text tokens to rank the significance of vision tokens. PyramidDrop~\cite{xing2024pyramiddrop} applies progressive pruning of tokens at different stages within the LLM.


\subsection{Video Visual Token Compression}

Nevertheless, when it comes to temporal dependencies between video frames, specialized compression designs need to be tailored. TempMe~\cite{shen2024tempme} extends progressive spatial tokens reduction by merging neighboring clips to minimize temporal redundancy.
DyCoke~\cite{tao2025dycoke} merges tokens across frames and applies dynamic KV cache reduction. However, its pruning during the prefilling stage struggles to achieve substantial token reduction while maintaining accuracy. PruneVID~\cite{huang2024prunevid} applies both merging and pruning across successive shallow LLM layers, but repeated pruning operations adversely affect overall efficiency. FastVID~\cite{shen2025fastvid} enhances compression by combining temporal segmentation with spatio-temporal token merging. FrameFusion~\cite{fu2024framefusion} applies both merging and pruning across successive shallow LLM layers, but repeated pruning operations adversely affect overall efficiency. Differently, our method proposes to construct compact token anchors across spatiotemporal to aggregate both local- and global-level semantic and contextual information through an optimization strategy.


